% Generated draft: requires mathematical and source review.
\documentclass[11pt]{article}
\usepackage[margin=25mm]{geometry}
\usepackage{fontspec}
\setmainfont{Latin Modern Roman}
\usepackage{amsmath,amssymb,mathtools,unicode-math}
\setmathfont{Latin Modern Math}
\usepackage{enumitem,needspace,xurl,hyperref}
\setlength{\parindent}{0pt}
\setlength{\parskip}{5pt}
\setlength{\emergencystretch}{3em}
\newcommand{\N}{\mathbb{N}}
\newcommand{\Z}{\mathbb{Z}}
\newcommand{\Q}{\mathbb{Q}}
\newcommand{\R}{\mathbb{R}}
\newcommand{\C}{\mathbb{C}}
\newcommand{\F}{\mathbb{F}}
\newcommand{\A}{\mathbb{A}}
\newcommand{\PP}{\mathbb{P}}
\newcommand{\E}{\mathbb{E}}
\newcommand{\eps}{\varepsilon}
\newcommand{\dd}{\,\mathrm{d}}
\begin{document}
\section*{UnsolvedMath ID: 4600009}
\textbf{Source problem code:} catalogue entry
\subsection{Definitions and mathematical statement}
A stationary Markov shift is a shift system \((X,\mu,\sigma)\) obtained from a stationary Markov chain on a specified alphabet, with \(X\) the allowed bi-infinite sequences and \((\sigma x)_i=x_{i+1}\). Mixing means
\[
\mu(A\cap\sigma^{-n}B)\longrightarrow\mu(A)\mu(B)
\]
for measurable \(A,B\). A finitary isomorphism is a measure-preserving isomorphism commuting with the shifts for which both coding directions have almost surely finite coding radii. A magic-word isomorphism is a further source-specific coding condition, not merely an arbitrary measurable or finitary isomorphism.

Write \(\beta_X\), \(\Delta_X\), \(\Gamma_X\), and \(g_X\in\Gamma_X/\Delta_X\) for the beta function, ratio group, weights group, and canonical quotient generator intended in the source. Their normalizations and the magic-word condition are not defined in the extracted record, so these names are retained as unresolved source-dependent invariants rather than replaced by invented formulas. Subject to those definitions, the stated conjecture is the equivalence
\[
X\text{ and }Y\text{ admit a magic-word isomorphism}
\quad\Longleftrightarrow\quad
\beta_X=\beta_Y,\quad\Delta_X=\Delta_Y,\quad g_X=g_Y,
\]
where equality of generators presupposes the intended identification of the quotient groups.

The source statement quantifies over two mixing Markov shifts without further qualifications. The supplied research summary instead describes pairs already finitarily isomorphic and satisfying a Markov-good condition; it does not define that condition. These are different hypothesis lists. Thus the equivalence above records the asserted target, but the quantification cannot be finalized until the admissible class, including finite versus countable alphabet, is checked. The summary calls the invariant equalities necessary and the sufficiency direction open through 2026, within its described setting; this does not establish the unrestricted version or a resolution.
\subsection{Short English statement}
Are the beta function, ratio group, and canonical weights-quotient generator complete invariants for magic-word isomorphism of the intended mixing Markov shifts? The precise invariant definitions and the extra admissibility hypotheses are missing.
\subsection{Sources}
\begin{itemize}
\item Local UnsolvedMath registry, permanent ID \(4600009\), exact source code \texttt{catalogue entry}, ``Good finitary conjecture.''
\item Boyle, \emph{Open Problems in Symbolic Dynamics} (2008), Problem/Question/Conjecture 9.1; \url{https://www.math.umd.edu/~mboyle/papers/openfinalsub3nov2007.pdf}. Recorded access 2026-07-29; pdf-text extraction.
\item Supplied research summary attributes necessary conditions to ``Keane--Smidt'' and the \(\Delta,\Gamma\) invariants to Boyle--Tuncel. It attributes the good finitary conjecture to Boyle, but supplies no complete references or definitions of the additional conditions.
\end{itemize}
\end{document}
